\documentclass[11pt]{article}
\usepackage[margin=2.0cm]{geometry}
\usepackage{booktabs}
\usepackage{amsmath}
\usepackage{amssymb}
\usepackage{enumitem}
\usepackage{array}
\usepackage{listings}
\usepackage{xcolor}
\usepackage[T1]{fontenc}
\usepackage{lmodern}

\setlength{\parindent}{0pt}
\setlength{\parskip}{3.6pt}
\newcommand{\code}[1]{\texttt{#1}}
\newcommand{\hr}{\vspace{2pt}\hrule\vspace{6pt}}
\renewcommand{\arraystretch}{1.12}

\definecolor{kw}{RGB}{0,90,160}
\definecolor{cmt}{RGB}{110,110,110}
\lstdefinestyle{vhdl}{
  language=VHDL, basicstyle=\ttfamily\footnotesize,
  keywordstyle=\color{kw}\bfseries, commentstyle=\color{cmt}\itshape,
  showstringspaces=false, columns=fullflexible, breaklines=true,
  frame=leftline, framesep=6pt, xleftmargin=8pt, aboveskip=6pt, belowskip=6pt
}
\lstset{style=vhdl}

\begin{document}

\begin{center}
{\Large\bfseries EE\,362 --- Computer Organization Laboratory}\\[2pt]
{\large\bfseries Lab 0 --- Student Practice Exercises: Solution}\\[2pt]
{\itshape P1 Voting Circuit \quad$\cdot$\quad P2 4-to-1 Multiplexer \quad$\cdot$\quad P3 3-to-8 Decoder}\\[4pt]
Fall 2026 \quad$\cdot$\quad Lab TA: Gandhi Torres
\end{center}
\hr

\textbf{Read this first.} Each of these three problems is a class example with one
more input. That is not a coincidence and it is not padding. In every case the
efficient solution is not a bigger truth table --- it is to \textbf{build the larger
circuit out of the smaller one you already have}.

\begin{center}
\begin{tabular}{llll}
\toprule
& class example & practice problem & built from \\
\midrule
P1 & majority of 3 & majority of 4 & \code{majority3} + \code{mux2} \\
P2 & 2-to-1 mux & 4-to-1 mux & three \code{mux2} \\
P3 & 2-to-4 decoder & 3-to-8 decoder & two \code{decoder2to4} \\
\bottomrule
\end{tabular}
\end{center}

That table is the method of the whole course in miniature. The processor you finish
in December is not a bigger truth table either --- it is seventeen components.

\vspace{4pt}\hr

\section*{P1 --- 4-Input Voting Circuit}

\subsection*{Specification}

Inputs $A,B,C,D$; output $Y = 1$ when \textbf{at least three} of the four are $1$.

With four voters, ``at least three'' is a strict majority --- so this is the
majority function again, one input wider.

\subsection*{Truth table --- Type E, 16 rows}

\begin{center}
\begin{tabular}{c|cccc|c|c}
\toprule
\# & $A$ & $B$ & $C$ & $D$ & ones & $Y$ \\
\midrule
0  & 0&0&0&0 & 0 & 0 \\
1  & 0&0&0&1 & 1 & 0 \\
2  & 0&0&1&0 & 1 & 0 \\
3  & 0&0&1&1 & 2 & 0 \\
4  & 0&1&0&0 & 1 & 0 \\
5  & 0&1&0&1 & 2 & 0 \\
6  & 0&1&1&0 & 2 & 0 \\
7  & 0&1&1&1 & 3 & \textbf{1} \\
8  & 1&0&0&0 & 1 & 0 \\
9  & 1&0&0&1 & 2 & 0 \\
10 & 1&0&1&0 & 2 & 0 \\
11 & 1&0&1&1 & 3 & \textbf{1} \\
12 & 1&1&0&0 & 2 & 0 \\
13 & 1&1&0&1 & 3 & \textbf{1} \\
14 & 1&1&1&0 & 3 & \textbf{1} \\
15 & 1&1&1&1 & 4 & \textbf{1} \\
\bottomrule
\end{tabular}
\end{center}

The ``ones'' column is not decoration --- it is the specification restated
numerically, and it makes the table self-checking. Add it to any table whose
specification is a counting condition.

Five of sixteen rows are $1$. Sanity check: that is $\binom{4}{3} + \binom{4}{4}
= 4 + 1 = 5$. \checkmark

\subsection*{Boolean derivation}

\textbf{(1) Canonical.}
\[
Y = \sum m(7,11,13,14,15) = \bar{A}BCD + A\bar{B}CD + AB\bar{C}D + ABC\bar{D} + ABCD
\]

\textbf{(2) Simplified.} Exactly the manoeuvre from Example 1, one input wider:
$m_{15} = ABCD$ is adjacent to each of the other four minterms, so by the
\textbf{idempotent law} write it four times:

\begin{align*}
Y &= \bar{A}BCD + A\bar{B}CD + AB\bar{C}D + ABC\bar{D}
     + \underbrace{ABCD + ABCD + ABCD + ABCD}_{\text{idempotent law}} \\
  &= (\bar{A}BCD + ABCD) + (A\bar{B}CD + ABCD) \\
  &\qquad + (AB\bar{C}D + ABCD) + (ABC\bar{D} + ABCD) \\
  &= BCD(\bar{A}+A) + ACD(\bar{B}+B) + ABD(\bar{C}+C) + ABC(\bar{D}+D) && \text{\small distributive law} \\[2pt]
Y &= \boxed{ABC + ABD + ACD + BCD} && \text{\small complement law}
\end{align*}

\textbf{This is minimal.} On a 4-variable K-map, each of $m_7, m_{11}, m_{13},
m_{14}$ is adjacent \emph{only} to $m_{15}$, so every group has size 2 and no
larger grouping exists. All four resulting prime implicants are essential.
Cost: 4 AND3 + 1 OR4.

\subsection*{The pattern --- worth more than the answer}

\begin{center}
\begin{tabular}{lll}
\toprule
function & result & terms \\
\midrule
at least 2 of 3 & $AB + AC + BC$ & $\binom{3}{2} = 3$ \\
at least 3 of 4 & $ABC + ABD + ACD + BCD$ & $\binom{4}{3} = 4$ \\
\bottomrule
\end{tabular}
\end{center}

\begin{quote}
\textbf{``At least $k$ of $n$'' is the OR of all $\binom{n}{k}$ products of $k$
variables.}
\end{quote}

Derive it once, then \emph{write down} the answer for any $n$ and $k$ without
touching a K-map. Note the cost, though: ``at least 5 of 9'' would be
$\binom{9}{5} = 126$ terms. Recognising that a formula explodes is itself a design
result --- at that size you stop using gates and build an adder tree instead.

\subsection*{Behavioural architecture}

\begin{lstlisting}
Y <= (A and B and C) or (A and B and D) or
     (A and C and D) or (B and C and D);
\end{lstlisting}

\subsection*{Structural architecture --- reuse both class components}

Split the problem on $D$'s vote:

\begin{itemize}[leftmargin=1.6em,itemsep=1pt]
\item If $D = 1$, we already have one yes --- we need only \textbf{2 more from
$A,B,C$}. That is \code{majority3}.
\item If $D = 0$, we need \textbf{all three} of $A,B,C$.
\end{itemize}

\[
Y = D \cdot \mathrm{Maj3}(A,B,C) \;+\; \bar{D} \cdot ABC
\]

Read the right-hand side: it is a \textbf{2-to-1 multiplexer selected by $D$}, with
\code{majority3} on one input and an AND3 on the other. Both parts came out of Lab 0
class examples, unmodified.

\begin{lstlisting}
architecture structural of voting4 is
    signal maj3, allABC : std_logic;
begin
    U_MAJ : entity work.majority3
        port map ( A => A, B => B, C => C, Y => maj3 );

    allABC <= A and B and C;

    U_MUX : entity work.mux2
        port map ( D0 => allABC, D1 => maj3, S => D, Y => Y );
end architecture structural;
\end{lstlisting}

\textbf{Proving the two agree.} Expand the structural form:
\begin{align*}
D(AB + AC + BC) + \bar{D}ABC
  &= ABD + ACD + BCD + ABC\bar{D} \\
\intertext{and expand the behavioural form's $ABC$ term by the complement law:}
ABC + ABD + ACD + BCD
  &= ABC(D + \bar{D}) + ABD + ACD + BCD \\
  &= \underbrace{ABCD}_{\text{absorbed into } ABD} + ABC\bar{D} + ABD + ACD + BCD \\
  &= ABD + ACD + BCD + ABC\bar{D} \qquad\checkmark && \text{\small absorption}
\end{align*}

Equal. \textbf{But the proof is not the verification} --- the testbench asserts
both architectures on all 16 vectors.

\subsection*{Testbench and result}

Four inputs, combinational: exhaustive, 16 vectors. Compute the expectation by
counting ones rather than typing a 16-entry literal table --- \textbf{a hand-typed
expected-value table can contain the same mistake as the design}.

\begin{lstlisting}
ones := 0;
for k in 0 to 3 loop
    if v(k) = '1' then ones := ones + 1; end if;
end loop;
if ones >= 3 then want := '1'; else want := '0'; end if;
\end{lstlisting}

Three checks per vector: behavioural vs.\ expected, structural vs.\ expected, and
the two against each other.

\begin{lstlisting}
# ** Note: PASS - voting4: 48 checks passed (16 vectors x 3).
\end{lstlisting}

\vspace{4pt}\hr

\section*{P2 --- 4-to-1 Multiplexer}

\subsection*{Truth table --- Type C}

Six inputs would give a 64-row exhaustive table. Do not write it. Write this:

\begin{center}
\begin{tabular}{c|cc|c}
\toprule
\# & $S_1$ & $S_0$ & $Y$ \\
\midrule
0 & 0 & 0 & $D_0$ \\
1 & 0 & 1 & $D_1$ \\
2 & 1 & 0 & $D_2$ \\
3 & 1 & 1 & $D_3$ \\
\bottomrule
\end{tabular}
\end{center}

Four rows for a six-input circuit, because the output column names a
\textbf{source} rather than a value. This is the form that survives when the mux is
32 bits wide in Lab 3.

\subsection*{Boolean derivation}

\[
Y = \bar{S_1}\bar{S_0}D_0 + \bar{S_1}S_0 D_1 + S_1\bar{S_0}D_2 + S_1 S_0 D_3
\qquad\text{or}\qquad
\boxed{Y = \sum_{k=0}^{3} m_k(S_1,S_0)\cdot D_k}
\]

Nothing simplifies, and nothing should --- but \emph{look at the coefficients}. The
four products $m_k(S_1,S_0)$ are precisely the four outputs of the 2-to-4 decoder
from class Example 3.

\begin{quote}
\textbf{A multiplexer is a decoder AND-ed with the data inputs, then OR-ed
together.} P2 and P3 are the same circuit seen from two directions: the decoder
produces the one-hot select, the mux consumes it.
\end{quote}

That relationship is why a large multiplexer and a large decoder cost roughly the
same, and it is worth one sentence in your report.

\subsection*{Two architectures, and a real trade-off}

\textbf{Behavioural --- flat select.} Two gate levels: AND3s in parallel, one OR4.

\begin{lstlisting}
S <= S1 & S0;
with S select
    Y <= D0   when "00",
         D1   when "01",
         D2   when "10",
         D3   when "11",
         'X'  when others;
\end{lstlisting}

\code{'X' when others} covers metavalues on the \emph{select}, not a fifth case ---
\code{"00"} through \code{"11"} already exhaust the two-bit space. Getting that
distinction right is the same trap that hid the enable bug in Example 3.

\textbf{Structural --- a mux tree.} $S_0$ chooses within each pair, $S_1$ chooses
between the pairs:

\begin{lstlisting}
U0 : entity work.mux2 port map ( D0 => D0, D1 => D1, S => S0, Y => L0 );
U1 : entity work.mux2 port map ( D0 => D2, D1 => D3, S => S0, Y => L1 );
U2 : entity work.mux2 port map ( D0 => L0, D1 => L1, S => S1, Y => Y  );
\end{lstlisting}

\begin{center}
\begin{tabular}{lll}
\toprule
& flat & tree \\
\midrule
depth & 2 gate levels & 2 mux levels ($\approx$ 4 gate levels) \\
fan-in & one OR4 & nothing wider than 2 \\
scaling to $2^n$ & one OR of width $2^n$ & $n$ levels, all narrow \\
\bottomrule
\end{tabular}
\end{center}

\textbf{Flat is faster here; the tree is what scales.} At 4 inputs the flat form
wins and you should use it. At 32 inputs the flat form needs a 32-input OR gate,
which no FPGA has --- the tool builds a tree anyway, and you have merely lost
control of how. Lab 3 is where this stops being hypothetical.

\subsection*{Testbench --- the important design decision}

64 vectors: \textbf{all 16 data patterns for each of the 4 select values.}

Sweeping every data pattern per select --- rather than just toggling the selected
input --- is what proves the mux \textbf{ignores the three inputs it did not
select}. That is the property most likely to be wrong and least likely to be tested.
A mux that ORs its inputs together passes any test where the unselected inputs
happen to be $0$.

\begin{lstlisting}
want := dv(s);      -- the selected input, whatever value it currently holds
\end{lstlisting}

One line, and it is the entire expected-value model. Compare it with typing a
64-row table.

\begin{lstlisting}
# ** Note: PASS - mux4: 192 checks passed (64 vectors x 3).
\end{lstlisting}

\vspace{4pt}\hr

\section*{P3 --- 3-to-8 Decoder with Enable}

\subsection*{Truth table --- Type E with a don't-care row}

\begin{center}
\begin{tabular}{c|cccc|cccccccc}
\toprule
\# & $E$ & $A_2$ & $A_1$ & $A_0$ & $D_7$ & $D_6$ & $D_5$ & $D_4$ & $D_3$ & $D_2$ & $D_1$ & $D_0$ \\
\midrule
0 & 0 & -- & -- & -- & 0&0&0&0&0&0&0&0 \\
1 & 1 & 0&0&0 & 0&0&0&0&0&0&0&\textbf{1} \\
2 & 1 & 0&0&1 & 0&0&0&0&0&0&\textbf{1}&0 \\
3 & 1 & 0&1&0 & 0&0&0&0&0&\textbf{1}&0&0 \\
4 & 1 & 0&1&1 & 0&0&0&0&\textbf{1}&0&0&0 \\
5 & 1 & 1&0&0 & 0&0&0&\textbf{1}&0&0&0&0 \\
6 & 1 & 1&0&1 & 0&0&\textbf{1}&0&0&0&0&0 \\
7 & 1 & 1&1&0 & 0&\textbf{1}&0&0&0&0&0&0 \\
8 & 1 & 1&1&1 & \textbf{1}&0&0&0&0&0&0&0 \\
\bottomrule
\end{tabular}
\end{center}

Nine rows, not sixteen: row 0's don't-cares collapse eight rows into one. That is a
\textbf{claim} --- that the address genuinely cannot reach the outputs when
disabled --- and the testbench has to prove it.

\subsection*{Boolean equations}

\[
D_k = E \cdot m_k(A_2,A_1,A_0), \qquad k = 0 \ldots 7
\]

Written out: $D_0 = E\bar{A_2}\bar{A_1}\bar{A_0}$, $D_1 = E\bar{A_2}\bar{A_1}A_0$,
\ldots, $D_7 = EA_2A_1A_0$. Eight equations, one per output; nothing simplifies,
because each is already a single product term.

\subsection*{Structural architecture --- the one that matters}

Build it from two \code{decoder2to4}. The key move:

\begin{quote}
\textbf{$A_2$ does not go into the sub-decoders. It splits the ENABLE.}
\end{quote}

\begin{lstlisting}
Elo <= E and (not A2);      -- lower half active when A2 = '0'
Ehi <= E and      A2;       -- upper half active when A2 = '1'

ULO : entity work.decoder2to4
    port map ( E => Elo, A1 => A1, A0 => A0, D => D(3 downto 0) );

UHI : entity work.decoder2to4
    port map ( E => Ehi, A1 => A1, A0 => A0, D => D(7 downto 4) );
\end{lstlisting}

Each sub-decoder sees only the low two address bits. The high bit decides
\emph{which sub-decoder is allowed to fire at all}, so exactly one output is high
across the pair --- one-hot is preserved by construction.

Check it against the table: $E=1$, $A_2A_1A_0 = 101$. Then \code{Ehi} = 1, the upper
decoder decodes $A_1A_0 = 01$ and drives its own output 1, which is \code{D(5)}.
Table row 6 says $D_5$. \checkmark

\textbf{Why this is the important one.} It composes without limit. A 4-to-16
decoder is two of these; a 5-to-32 decoder is two 4-to-16s. \textbf{5-to-32 is
exactly what the register file needs in Lab 6} to turn a 5-bit register number into
32 write-enables. You have now built the mechanism; Lab 6 only widens it.

\subsection*{Testbench --- and the bug it is aimed at}

16 vectors: 8 addresses $\times$ $E \in \{0,1\}$. \textbf{The expectation is
computed, not typed:}

\begin{lstlisting}
if ev = 0 then
    want := (others => '0');
else
    want := std_logic_vector(shift_left(to_unsigned(1, 8), a));
end if;
\end{lstlisting}

``One, shifted left by the address'' \emph{is} the definition of a decoder, stated
in one line. It cannot drift out of step with the specification the way a typed
table can.

Plus a \textbf{one-hot invariant} --- count the set bits and require exactly one
when enabled, zero when not. An invariant catches faults a value comparison can
miss, and it is the first time in this course you will write a check that is not a
lookup.

\textbf{And the point of the whole exercise:} the eight vectors with $E = 0$. The
behavioural decoder shown in class ignored its enable entirely, and \emph{a
testbench that only drives $E = 1$ passes it}. Row 0 of the table made a claim; these
eight vectors are what turn that claim into evidence.

\begin{lstlisting}
# ** Note: PASS - decoder3to8: 48 checks passed (16 vectors x 3).
\end{lstlisting}

\vspace{4pt}\hr

\section*{Quartus results --- interpreting all three}

Synthesised for Cyclone V, device \code{5CSEMA4U23C6}:

\begin{center}
\begin{tabular}{lccc}
\toprule
& \code{voting4} & \code{mux4} & \code{decoder3to8} \\
\midrule
Total pins & 5 & 7 & 12 \\
\quad in / out & 4 / 1 & 6 / 1 & 4 / 8 \\
Logic utilisation (ALMs) & 1 & 2 & 5 \\
\textbf{Total registers} & \textbf{0} & \textbf{0} & \textbf{0} \\
Block memory bits & 0 & 0 & 0 \\
\bottomrule
\end{tabular}
\end{center}

All three are under 1\% of the device's 15{,}880 ALMs. The absolute numbers are
uninteresting; the \emph{ordering} is not. \code{voting4} fits in a
\textbf{single ALM} --- it is one function of four inputs, and a Cyclone V adaptive
logic module can implement any function of up to six. All that Boolean simplification
work bought exactly nothing in hardware here, because the tool would have fitted the
canonical five-minterm form into the same one ALM. \textbf{Simplification is for your
understanding and your report; the fitter does its own.} It starts to matter for real
in Lab 5, where the ALU is too big to absorb.

\code{decoder3to8} costs the most (5 ALMs) despite having the fewest distinct input
combinations, because \textbf{ALM cost tracks outputs, not inputs} --- eight outputs
need eight functions, and one ALM can produce at most two.

\textbf{Zero registers, in all three.} All three circuits are purely combinational
--- output depends only on present inputs, nothing is stored, and there is no clock
anywhere. If your Flow Summary shows a register or a latch, you have accidentally
described memory: check for a \code{process} with an incomplete \code{if} or
\code{case}. That is the Lab 1 mechanism appearing where you did not want it.

\textbf{Pin counts are the shape of the circuit.} \code{decoder3to8} is the only one
of the three with more outputs than inputs --- 4 in, 8 out --- because a decoder's
job is precisely to expand a compact code into one wire per case. \code{mux4} is the
opposite: 6 in, 1 out, collapsing many wires to one. \textbf{Muxes narrow, decoders
widen}, and you will use them in exactly those roles for the rest of the semester.

\textbf{Timing Analyzer: \code{No clocks defined in design}, no $F_{max}$ --- and
that is correct.}
$F_{max}$ answers ``how fast can I clock this?'' and none of these has a clock.
Report the combinational propagation delay if you want a number, but do not report a
missing $F_{max}$ as a failure. The first real $F_{max}$ of the course arrives in
Lab 2.

\textbf{RTL Viewer.} Compare the two architectures of each design deliberately ---
this is the most instructive ten minutes in the whole practice set:

\begin{itemize}[leftmargin=1.6em,itemsep=1pt]
\item \code{voting4} --- behavioural shows four AND3s into an OR4; structural shows
the \code{majority3} block and a mux. \emph{Different pictures, same function.}
\item \code{mux4} --- flat vs.\ the two-level tree. You can \textbf{see} the depth
difference that the trade-off table describes.
\item \code{decoder3to8} --- structural shows two \code{decoder2to4} blocks with the
enable split. Confirm $A_2$ really does not enter either sub-block.
\end{itemize}

Then note what Quartus does about it: after optimisation the two architectures
usually map to \emph{identical} logic. \textbf{Your structural choice is about how
the design is written, understood, and reused --- not about what reaches the
silicon.} That is the honest lesson, and it is the reason the rest of this course
builds structurally even when behavioural would fit.

\section*{Files}

All under \code{Lab\_0/practice/solution/}, all verified in Questa:

\begin{center}
\begin{tabular}{llc}
\toprule
design & testbench & result \\
\midrule
\code{voting4.vhd} (2 arch.) & \code{tb\_voting4.vhd} & PASS 48 \\
\code{mux4.vhd} (2 arch.) & \code{tb\_mux4.vhd} & PASS 192 \\
\code{decoder3to8.vhd} (2 arch.) & \code{tb\_decoder3to8.vhd} & PASS 48 \\
\midrule
\multicolumn{2}{l}{shared: \code{majority3.vhd}, \code{mux2.vhd}, \code{decoder2to4.vhd}} & \textbf{288} \\
\bottomrule
\end{tabular}
\end{center}

\code{decoder2to4.vhd} here is the \textbf{corrected} version --- it gates on
\code{E}. The file shown in class did not.

\vspace{6pt}
\hr
{\small \textbf{One thing to take from all three.} Not one of these problems was
solved by writing a bigger truth table. P1 split on one input and reused a
3-input component; P2 was a tree of 2-input components; P3 was two half-sized
decoders with the enable split. \textbf{Composition, not enumeration} --- that is
the only technique that still works in December, when the thing being built is a
processor.}

\end{document}
